%% toc: P1 3 3.1,3.2,3.6
\section{Proofs for orders, natural numbers, and cardinal assignment}
\label{sec:annex-p01-w02-p01-008}

The ambient theory is the classical set theory of the metalanguage fixed at
the start of \cref{ch:set-prerequisites}. The Axiom of Choice is used only in
the proofs that name it. ``Ordinal item~\(k\)'' refers to item~\(k\) of
\cref{thm:p01-w02-p01-011-ordinals}: every member of an ordinal is an
ordinal and no ordinal is a member of itself (item~1); \(\alpha\subseteq\beta\)
implies \(\alpha=\beta\) or \(\alpha\in\beta\) (item~2); exactly one of
\(\alpha\in\beta\), \(\alpha=\beta\), \(\beta\in\alpha\) holds, and
\(\alpha\leq\beta\) iff \(\alpha\subseteq\beta\) (item~3); every nonempty set
of ordinals has a unique least member (item~4); and every transitive set of
ordinals is an ordinal, in particular \(\alpha+1\) (item~5). Items~1--5 are
proved there from \cref{def:ordinal-cardinal} without natural numbers,
induction or recursion.

\begin{annexitem}{annex:p01-w02-p01-008-strict-nonstrict}{thm:p01-w02-p01-008-strict-nonstrict}{Strict and non-strict orders correspond}
\begin{proof}
\emph{Item~1.} Let \(R\) be irreflexive and transitive on \(A\). If \(aRb\)
and \(bRa\), transitivity gives \(aRa\), against irreflexivity; so \(R\) is
asymmetric. Put \(Q=R\cup\mathrm{id}_A\subseteq A\times A\). It is reflexive
because \(\mathrm{id}_A\subseteq Q\). It is antisymmetric: if \(aQb\),
\(bQa\) and \(a\neq b\), neither pair lies in \(\mathrm{id}_A\), so \(aRb\)
and \(bRa\), against asymmetry. It is transitive: let \(aQb\) and \(bQc\).
If \(a=b\), then \(aQc\) is the second hypothesis; if \(b=c\), it is the
first. Otherwise \(aRb\) and \(bRc\), so \(aRc\) and \(aQc\). Finally,
irreflexivity says \(R\cap\mathrm{id}_A=\varnothing\), so
\(Q\setminus\mathrm{id}_A=R\setminus\mathrm{id}_A=R\).

\emph{Item~2.} Let \(Q\) be a partial order on \(A\) and
\(P=Q\setminus\mathrm{id}_A\). Then \(P\) is irreflexive by construction.
Let \(aPb\) and \(bPc\). Then \(aQb\) and \(bQc\), so \(aQc\). If \(a=c\),
then \(aQb\) and \(bQa\), so \(a=b\) by antisymmetry, against \(aPb\).
Hence \(a\neq c\) and \(aPc\). Since \(Q\) is reflexive,
\(\mathrm{id}_A\subseteq Q\), so \(P\cup\mathrm{id}_A=Q\).

\emph{The bijection.} Items~1 and~2 show that each map sends orders of one
kind to orders of the other kind and that each composite is the identity, in
the displayed equations \((R\cup\mathrm{id}_A)\setminus\mathrm{id}_A=R\) and
\((Q\setminus\mathrm{id}_A)\cup\mathrm{id}_A=Q\). Two maps with both
composites the identity are mutually inverse bijections.

\emph{Item~3.} Let \(R\) be a strict partial order and
\(Q=R\cup\mathrm{id}_A\), a partial order by item~1. Suppose \(R\) is a
strict linear order and let \(a,b\in A\). If \(a=b\), then \(aQa\). If
\(a\neq b\), then \(aRb\) or \(bRa\), so \(aQb\) or \(bQa\). Hence \(Q\) is
a non-strict linear order. Conversely suppose \(Q\) is a non-strict linear
order and let \(a\neq b\). Then \(aQb\) or \(bQa\), and since the pair is
not in \(\mathrm{id}_A\), \(aRb\) or \(bRa\). Both cannot hold, by
asymmetry (item~1). With irreflexivity and transitivity, \(R\) is a strict
linear order.

\emph{Item~4.} Let \(R\) be a strict linear order, \(B\subseteq A\) and
\(b\in B\). Suppose \(b\) is \(R\)-minimal in \(B\), and let \(a\in B\). If
\(a=b\), then \(bQa\). If \(a\neq b\), then \(aRb\) or \(bRa\); the first is
excluded by minimality, so \(bRa\) and \(bQa\). Hence \(b\) is \(Q\)-least.
Conversely suppose \(bQa\) for every \(a\in B\), and suppose some \(a\in B\)
satisfies \(aRb\). Then \(a\neq b\) by irreflexivity, and \(bQa\) with
\(a\neq b\) gives \(bRa\), against asymmetry. Hence \(b\) is \(R\)-minimal.

\emph{Item~5.} By item~3, \(R\) is a strict linear order iff \(Q\) is a
non-strict linear order; when they are not, neither side of item~5 holds.
When they are, item~4 shows that a nonempty
\(B\subseteq A\) has an \(R\)-minimal member iff it has a \(Q\)-least
member, and the member is the same. So every nonempty subset has an
\(R\)-minimal member iff every nonempty subset has a \(Q\)-least member,
which is the claim.
\end{proof}
\end{annexitem}

\begin{annexitem}{annex:p01-w02-p01-008-omega}{thm:p01-w02-p01-008-omega}{Existence and leastness of \(\omega\)}
\begin{proof}
The set \(\omega\) exists by separation from \(I_0\) with the formula
``\(x\) belongs to every inductive set'', in which the inductive condition
is written out as in \cref{def:p01-w02-p01-008-inductive}.

\emph{\(\omega\) is inductive.} The empty set belongs to every inductive set
by definition, and \(\varnothing\in I_0\); so \(\varnothing\in\omega\). Let
\(x\in\omega\). Then \(x\) belongs to every inductive set \(J\), so
\(S(x)\in J\) for every inductive \(J\); in particular \(S(x)\in I_0\).
Hence \(S(x)\in\omega\).

\emph{\(\omega\) is included in every inductive set.} If \(x\in\omega\) and
\(J\) is inductive, then \(x\in J\) by the defining condition.

\emph{Independence of \(I_0\).} Let \(I_1\) be another inductive set and
\(\omega_1\) the set formed from it in the same way. If \(x\in\omega\), then
\(x\) belongs to every inductive set, so \(x\in I_1\) and \(x\in\omega_1\);
symmetrically \(\omega_1\subseteq\omega\). By Extensionality
\(\omega=\omega_1\).

\emph{Uniqueness.} If \(\omega'\) is inductive and included in every
inductive set, then \(\omega'\subseteq\omega\) because \(\omega\) is
inductive, and \(\omega\subseteq\omega'\) because \(\omega'\) is inductive.
So \(\omega'=\omega\).
\end{proof}
\end{annexitem}

\begin{annexitem}{annex:p01-w02-p01-008-induction}{thm:p01-w02-p01-008-induction}{Induction on \(\N\)}
\begin{proof}
\emph{Item~1.} The hypotheses say that \(P\) is inductive: \(0=\varnothing\in P\),
and if \(n\in P\) then \(S(n)=n+1\in P\). By
\cref{thm:p01-w02-p01-008-omega}, \(\N=\omega\subseteq P\). With
\(P\subseteq\N\) this gives \(P=\N\).

\emph{Item~2.} By separation, \(P=\{n\in\N:\varphi(n)\}\) is a set, with
the same parameters as \(\varphi\). The hypotheses give \(0\in P\) and
\(n+1\in P\) whenever \(n\in P\). By item~1, \(P=\N\), that is,
\(\varphi(n)\) for every \(n\in\N\).
\end{proof}
\end{annexitem}

\begin{annexitem}{annex:p01-w02-p01-008-naturals}{thm:p01-w02-p01-008-naturals}{Order structure of the natural numbers}
\begin{proof}
\emph{Item~1.} Let \(P\) be the set of \(n\in\N\) such that \(n\) is an
ordinal and \(n\subseteq\N\). The empty set is transitive, and \(\in\)
strictly well-orders it because every clause of that definition is
vacuous; and \(\varnothing\subseteq\N\). So \(0\in P\). Let \(n\in P\). Then
\(n+1=n\cup\{n\}\) is an ordinal by ordinal item~5, and
\(n+1\subseteq\N\) because \(n\subseteq\N\) and \(n\in\N\). So \(n+1\in P\).
By \cref{thm:p01-w02-p01-008-induction}, \(P=\N\). Thus every natural
number is an ordinal, and every member of a natural number is a natural
number, which says that \(\N\) is transitive.

\emph{Item~2.} By item~1, \(\N\) is a transitive set of ordinals, so it is
an ordinal by ordinal item~5, and \(\in\) strictly well-orders \(\N\). For
\(m,n\in\N\): if \(m\neq n\), exactly one of \(m\in n\) and \(n\in m\)
holds; if \(m=n\), neither holds, by irreflexivity. Every nonempty subset
has an \(\in\)-minimal member, which by item~4 of
\cref{thm:p01-w02-p01-008-strict-nonstrict} is a least member for
\(\leq\).

\emph{Item~3.} By item~1, \(n\subseteq\N\), so
\(n=\{k\in\N:k\in n\}=\{k\in\N:k<n\}\). Since \(n+1=n\cup\{n\}\), a natural
number \(k\) lies in \(n+1\) iff \(k<n\) or \(k=n\), that is, iff
\(k\leq n\). Hence \(0=\{k\in\N:k<0\}=\varnothing\), \(1=0+1=\{0\}\), and
\(2=1+1=\{k\in\N:k\leq1\}=\{0,1\}\). For \(n\neq0\), item~4 gives
\(n=(n-1)+1=\{k\in\N:k\leq n-1\}\): the members of \(n\) are exactly the
natural numbers from \(0\) up to \(n-1\) in the order \(<\). If \(k\in\N\)
and \(k<n+1\), then \(k\leq n\), so \(n<k\) fails by item~2; no natural
number lies strictly between \(n\) and \(n+1\).

\emph{Item~4.} Since \(n\in n+1\) and \(0\) has no members, \(0\neq n+1\).
Let \(P\) be the set of \(n\in\N\) such that \(n=0\) or \(n=m+1\) for some
\(m\in\N\). Then \(0\in P\), and for \(n\in\N\), \(n+1\in P\) with
witness \(m=n\). By \cref{thm:p01-w02-p01-008-induction}, \(P=\N\). For
uniqueness let \(m+1=m'+1\) with \(m,m'\in\N\) and suppose \(m\neq m'\).
From \(m\in m'+1\) and \(m\neq m'\) we get \(m\in m'\); likewise
\(m'\in m\). This contradicts item~2. So \(m=m'\), which is also the last
claim of the item.

\emph{Item~5.} \(\omega\neq\varnothing\) because \(0\in\omega\). Suppose
\(\omega=S(x)\) for a set \(x\). Then \(x\in S(x)=\omega\), and since
\(\omega\) is inductive, \(S(x)\in\omega\), that is, \(\omega\in\omega\).
This contradicts ordinal item~1, as \(\omega\) is an ordinal by item~2. In
particular \(\omega\neq\alpha+1\) for every ordinal \(\alpha\), so \(\omega\)
is a limit ordinal as defined after \cref{def:ordinal-cardinal}.
\end{proof}
\end{annexitem}

\begin{annexitem}{annex:p01-w02-p01-008-pigeonhole}{thm:p01-w02-p01-008-pigeonhole}{Pigeonhole principle}
\begin{proof}
\emph{Item~1.} By induction on \(n\)
(\cref{thm:p01-w02-p01-008-induction}, item~2). For \(n=0\): \(1=\{0\}\)
is nonempty and \(0=\varnothing\), and no function from a nonempty set to
\(\varnothing\) exists. Assume that no injection \(n+1\to n\) exists, and
suppose \(f:(n+1)+1\to n+1\) is injective. Write \(s=n+1\). Then
\((n+1)+1=s\cup\{s\}\), and \(s\notin s\) by ordinal item~1, so
\(s\notin n+1\).

\emph{Case A: \(n\notin f[n+1]\).} For \(x\in n+1\), \(f(x)\in n+1=n\cup\{n\}\)
and \(f(x)\neq n\), so \(f(x)\in n\). The restriction of \(f\) to \(n+1\) is
an injection \(n+1\to n\), against the assumption.

\emph{Case B: \(n=f(k)\) for some \(k\in n+1\).} As \(k\in n+1\) and
\(s\notin n+1\), \(s\neq k\), so \(f(s)\neq f(k)=n\) and \(f(s)\in n\).
Define \(g\) on \(n+1\) by \(g(k)=f(s)\) and \(g(x)=f(x)\) for
\(x\neq k\). For \(x\neq k\), \(f(x)\neq f(k)=n\), so \(f(x)\in n\); thus
\(g\) maps \(n+1\) into \(n\). Let \(x\neq y\) in \(n+1\). If neither is
\(k\), \(g(x)=f(x)\neq f(y)=g(y)\). If \(x=k\neq y\), then \(g(x)=g(y)\)
would mean \(f(s)=f(y)\), hence \(s=y\in n+1\), which is false. So \(g\) is
an injection \(n+1\to n\), against the assumption. Both cases are
impossible, so no injection \((n+1)+1\to n+1\) exists.

\emph{Item~2.} Let \(m<n\). By item~3 of \cref{thm:p01-w02-p01-008-naturals},
\(m+1=\{k\in\N:k\leq m\}\), and \(k\leq m<n\) gives \(k<n\) by item~2 there;
so \(m+1\subseteq n\). An injection \(n\to m\) would restrict to an injection
\(m+1\to m\), against item~1. If \(m\neq n\), then \(m<n\) or \(n<m\); a
bijection between \(m\) and \(n\), or its inverse, would be an injection from
the larger into the smaller.

\emph{Item~3.} An ordinal smaller than \(n\) is a member \(k\in n\); then
\(k\in\N\) and \(k<n\). By item~2, \(k\) is not bijective with \(n\). So
\(n\) is a cardinal.

\emph{Item~4.} Since \(\N\) is transitive and \(n\in\N\), \(n+1\subseteq\N\).
An injection \(\N\to n\) would restrict to an injection \(n+1\to n\),
against item~1. If \(\N\) were finite, a bijection \(n\to\N\) would have an
inverse that is an injection \(\N\to n\). The ordinals smaller than
\(\N=\omega\) are its members \(n\), and a bijection between \(n\) and \(\N\)
would again give an injection \(\N\to n\). So \(\N\) is a cardinal, and it is
bijective with no \(n\in\N\), which is the definition of an infinite cardinal.

\emph{Item~5.} If \(f:n\to A\) and \(g:m\to A\) are bijections, then
\(g^{-1}\circ f\) is a bijection \(n\to m\), so \(n=m\) by item~2.

\emph{Item~6.} If \(f:n\to A\) is a bijection, then \(f^{-1}:A\to n\) is an
injection with values in \(n\subseteq\N\), so \(A\) is countable. Now let
\(f:A\to\N\) be an injection and \(X=f[A]\); then \(f\) is a bijection
\(A\to X\). The set \(X\subseteq\omega\) is strictly well-ordered by
\(\in\), as a subset of an ordinal. By
\cref{thm:p01-w02-p01-008-representation}, item~1, there are an ordinal
\(\alpha\) and an order isomorphism \((X,\in)\cong(\alpha,\in)\), and by
item~2 there \(\alpha\leq\omega\). If \(\alpha=\omega\), then \(A\) is
bijective with \(X\) and so with \(\N\): \(A\) is countably infinite. If
\(\alpha\in\omega\), then \(\alpha=n\) for some \(n\in\N\), and \(A\) is
bijective with \(n\): \(A\) is finite. A set that is both finite and
countably infinite would give a bijection between some \(n\in\N\) and
\(\N\), against item~4.

\emph{Item~7.} Every \(n\in\N\) is an ordinal
(\cref{thm:p01-w02-p01-008-naturals}, item~1) and is finite by the identity
bijection \(n\to n\). Conversely let \(\alpha\) be an ordinal and
\(f:n\to\alpha\) a bijection with \(n\in\N\). By ordinal item~3,
\(\alpha\in\omega\), \(\alpha=\omega\) or \(\omega\in\alpha\). In the last
two cases \(\omega\subseteq\alpha\), because \(\alpha\) is transitive, and
\(f^{-1}\) restricted to \(\omega\) is an injection \(\N\to n\), against
item~4. So \(\alpha\in\omega\), that is, \(\alpha\) is a natural number.
\end{proof}
\end{annexitem}

\begin{annexitem}{annex:p01-w02-p01-008-representation}{thm:p01-w02-p01-008-representation}{Every well-order is isomorphic to exactly one ordinal}
\begin{proof}
\emph{Item~1, existence.} By separation let \(r=\{x\in W:x\notin x\}\).
Then \(r\notin W\): if \(r\in W\), the defining condition gives \(r\in r\)
iff \(r\notin r\). Put \(W'=W\cup\{r\}\) and
\(\prec'=\prec\cup\{(w,r):w\in W\}\), a subset of \(W'\times W'\). We check
that \(\prec'\) strictly well-orders \(W'\). The element \(r\) never stands
on the left of \(\prec'\), because the left entries of \(\prec\) and of the
added pairs lie in \(W\); so \(r\prec'r\) fails, and \(\prec'\) is
irreflexive. If \(x\prec'y\prec'z\), then \(x\in W\); if \(z\in W\), the
pair \((y,z)\) lies in \(\prec\), so \(y\in W\) and \(x\prec z\); if
\(z=r\), then \(x\prec'r\) by definition. For distinct \(x,y\in W\) exactly
one of \(x\prec'y\), \(y\prec'x\) holds because this is so for \(\prec\); for
\(w\in W\), \(w\prec'r\) holds and \(r\prec'w\) fails. Let
\(\varnothing\neq B\subseteq W'\). If \(B\cap W\neq\varnothing\), its
\(\prec\)-minimal member \(b\) is \(\prec'\)-minimal in \(B\), because
\(r\prec'b\) fails and \(a\prec'b\) iff \(a\prec b\) for \(a\in W\).
Otherwise \(B=\{r\}\) and \(r\) is minimal. In \((W',\prec')\) the initial
segment below \(r\) is \(\{w\in W':w\prec'r\}=W\), ordered by \(\prec\). By
item~2 of \cref{thm:p01-w02-p01-011-representation}, applied to \(W'\) and
\(r\), it is order isomorphic to an ordinal.

\emph{Item~1, uniqueness.} Let \(f:(W,\prec)\cong(\alpha,\in)\) and
\(g:(W,\prec)\cong(\alpha',\in)\). Then \(g\circ f^{-1}\) is an order
isomorphism \(\alpha\to\alpha'\). By item~1 of
\cref{thm:p01-w02-p01-011-representation}, \(\alpha=\alpha'\) and
\(g\circ f^{-1}\) is the identity of \(\alpha\). Hence \(g=f\).

\emph{Item~2.} Let \(h:(X,\in)\cong(\alpha,\in)\) and
\(g=h^{-1}:\alpha\to X\). Members of \(X\) are members of \(\beta\), hence
ordinals (ordinal item~1), and \(\gamma\in\delta\) iff
\(g(\gamma)\in g(\delta)\) for \(\gamma,\delta\in\alpha\). We claim
\(\gamma\leq g(\gamma)\) for every \(\gamma\in\alpha\). Otherwise, by
ordinal item~3, the set \(D=\{\gamma\in\alpha:g(\gamma)\in\gamma\}\) is
nonempty; let \(\gamma\) be its least member (ordinal item~4) and
\(\eta=g(\gamma)\). Then \(\eta\in\gamma\), so \(\eta\in\alpha\) because
\(\alpha\) is transitive, and \(\eta\in\gamma\) gives
\(g(\eta)\in g(\gamma)=\eta\). Thus \(\eta\in D\) and \(\eta\in\gamma\),
against the leastness of \(\gamma\) (ordinal item~3 excludes
\(\gamma\leq\eta\) together with \(\eta\in\gamma\)). Hence
\(\gamma\leq g(\gamma)\) for every \(\gamma\in\alpha\). As
\(g(\gamma)\in\beta\) and \(\beta\) is transitive, \(\gamma\in\beta\). So
\(\alpha\subseteq\beta\), and \(\alpha\leq\beta\) by ordinal item~3.
\end{proof}
\end{annexitem}

\begin{annexitem}{annex:p01-w02-p01-008-well-ordering}{thm:p01-w02-p01-008-well-ordering}{Well-ordering}
\begin{proof}
Let \(A\) be a set and \(X=\mathcal P(A)\setminus\{\varnothing\}\), a set
of nonempty sets with \(\bigcup X\subseteq A\). The Axiom of Choice gives a
function \(c:X\to\bigcup X\) with \(c(Y)\in Y\) for every \(Y\in X\). This
is the only use of the axiom.

\emph{The successor operation on subsets.} For \(B\subseteq A\) with
\(B\neq A\), the set \(A\setminus B\) is a nonempty subset of \(A\), so
\(c(A\setminus B)\) is defined and lies in \(A\setminus B\); put
\(B^{+}=B\cup\{c(A\setminus B)\}\). Put \(A^{+}=A\). In both cases
\(B\subseteq B^{+}\subseteq A\), and if \(B\neq A\) then
\(B^{+}\setminus B=\{c(A\setminus B)\}\), a single element.

\emph{Towers.} A family \(\mathcal T\subseteq\mathcal P(A)\) is a tower iff
(T1) \(\varnothing\in\mathcal T\); (T2) \(B\in\mathcal T\) implies
\(B^{+}\in\mathcal T\); and (T3) \(\bigcup\mathcal S\in\mathcal T\) for every
\(\mathcal S\subseteq\mathcal T\). The family \(\mathcal P(A)\) is a tower.
By separation from \(\mathcal P(\mathcal P(A))\) let \(\mathcal T_0\) be the
set of \(B\subseteq A\) that belong to every tower. Then \(\mathcal T_0\)
is a tower: \(\varnothing\) belongs to every tower; if \(B\) belongs to every
tower, so does \(B^{+}\); and if \(\mathcal S\subseteq\mathcal T_0\), then
\(\mathcal S\) is included in every tower \(\mathcal T\), so
\(\bigcup\mathcal S\in\mathcal T\) for every tower. Also
\(\mathcal T_0\subseteq\mathcal T\) for every tower \(\mathcal T\). Hence
\emph{every tower \(\mathcal U\subseteq\mathcal T_0\) equals
\(\mathcal T_0\).}

\emph{Step~1: a property of comparable sets.} Call \(C\in\mathcal T_0\)
comparable iff for every \(B\in\mathcal T_0\), \(B\subseteq C\) or
\(C\subseteq B\). We show: if \(C\) is comparable, then every
\(B\in\mathcal T_0\) satisfies \(B\subseteq C\) or \(C^{+}\subseteq B\). Let
\(\mathcal U_C\) be the set of \(B\in\mathcal T_0\) with this property; we
show it is a tower. (T1): \(\varnothing\subseteq C\). (T2): let
\(B\in\mathcal U_C\). If \(C^{+}\subseteq B\), then
\(C^{+}\subseteq B\subseteq B^{+}\). If \(B=C\), then \(C^{+}=B^{+}\). If
\(B\subsetneq C\), then \(B\neq A\), and since \(C\) is comparable,
\(B^{+}\subseteq C\) or \(C\subseteq B^{+}\). In the second case
\(C\setminus B\) is nonempty and included in \(B^{+}\setminus B\), which has
exactly one element; so \(C\setminus B=B^{+}\setminus B\) and \(C=B^{+}\).
Either way \(B^{+}\subseteq C\). So \(B^{+}\in\mathcal U_C\). (T3): let
\(\mathcal S\subseteq\mathcal U_C\). If every member of \(\mathcal S\) is
included in \(C\), so is \(\bigcup\mathcal S\); otherwise some member
includes \(C^{+}\), and so does \(\bigcup\mathcal S\). Hence \(\mathcal U_C\)
is a tower included in \(\mathcal T_0\), so \(\mathcal U_C=\mathcal T_0\).

\emph{Step~2: \(\mathcal T_0\) is linearly ordered by inclusion.} Let
\(\mathcal K\) be the set of comparable members of \(\mathcal T_0\); we show
it is a tower. (T1): \(\varnothing\subseteq B\) for every \(B\). (T2): let
\(C\in\mathcal K\) and \(B\in\mathcal T_0\). By step~1, \(B\subseteq C\subseteq
C^{+}\) or \(C^{+}\subseteq B\); so \(C^{+}\) is comparable. (T3): let
\(\mathcal S\subseteq\mathcal K\) and \(B\in\mathcal T_0\). If
\(B\subseteq S\) for some \(S\in\mathcal S\), then
\(B\subseteq\bigcup\mathcal S\); otherwise \(S\subseteq B\) for every
\(S\in\mathcal S\), and \(\bigcup\mathcal S\subseteq B\). So
\(\mathcal K=\mathcal T_0\): any two members of \(\mathcal T_0\) are
comparable under \(\subseteq\), and step~1 applies to every
\(C\in\mathcal T_0\).

\emph{Step~3: the sets \(T_x\).} For \(x\in A\) let
\(T_x=\bigcup\{B\in\mathcal T_0:x\notin B\}\). By (T3),
\(T_x\in\mathcal T_0\); \(x\notin T_x\); and every \(B\in\mathcal T_0\) with
\(x\notin B\) satisfies \(B\subseteq T_x\). As \(x\in A\setminus T_x\),
\(T_x\neq A\), so \(T_x^{+}=T_x\cup\{c(A\setminus T_x)\}\in\mathcal T_0\)
strictly includes \(T_x\). If \(x\notin T_x^{+}\), then
\(T_x^{+}\subseteq T_x\), which is false. So
\(x\in T_x^{+}\setminus T_x=\{c(A\setminus T_x)\}\), that is,
\(x=c(A\setminus T_x)\). Consequently \(T_x=T_y\) implies \(x=y\).

\emph{Step~4: the order.} By separation from \(A\times A\), let \(x\prec y\)
iff \(x\in T_y\). We show that \(x\prec y\) iff \(T_x\subsetneq T_y\). If
\(x\in T_y\), then \(T_y\subseteq T_x\) is impossible because
\(x\notin T_x\); by step~2, \(T_x\subseteq T_y\), and \(x\in T_y\setminus T_x\)
makes the inclusion strict. If \(T_x\subsetneq T_y\), step~1 with \(C=T_x\)
and \(B=T_y\) gives \(T_y\subseteq T_x\) or \(T_x^{+}\subseteq T_y\); the
first is false, so \(x\in T_x^{+}\subseteq T_y\).

\emph{Step~5: \(\prec\) is a strict linear order.} It is irreflexive
because \(x\notin T_x\), and transitive because strict inclusion is. For
distinct \(x,y\in A\), \(T_x\neq T_y\) by step~3, and by step~2 one of
\(T_x\subsetneq T_y\), \(T_y\subsetneq T_x\) holds; both cannot hold, since
together they give \(T_x\subsetneq T_x\). By step~4, exactly one of
\(x\prec y\), \(y\prec x\) holds.

\emph{Step~6: minimal members.} Let \(\varnothing\neq D\subseteq A\) and
\(U=\bigcup\{B\in\mathcal T_0:B\cap D=\varnothing\}\). Then
\(U\in\mathcal T_0\) and \(U\cap D=\varnothing\), so \(U\neq A\) because
\(D\) is nonempty. Hence \(U^{+}=U\cup\{d\}\) with \(d=c(A\setminus U)\),
and \(U^{+}\in\mathcal T_0\) strictly includes \(U\). If
\(U^{+}\cap D=\varnothing\), then \(U^{+}\subseteq U\), which is false;
since \(U\cap D=\varnothing\), the common element is \(d\), so \(d\in D\).
We show \(T_d=U\). As \(d\notin U\), \(U\subseteq T_d\). By step~1 with
\(C=U\) and \(B=T_d\), \(T_d\subseteq U\) or \(U^{+}\subseteq T_d\); the
second would give \(d\in T_d\), which is false. So \(T_d=U\). Now if
\(e\in D\), then \(e\notin U=T_d\), so \(e\prec d\) fails. Thus \(d\) is
\(\prec\)-minimal in \(D\).

By steps~5 and~6, \(\prec\) strictly well-orders \(A\).
\end{proof}
\end{annexitem}

\begin{annexitem}{annex:p01-w02-p01-008-cardinal-assignment}{thm:p01-w02-p01-008-cardinal-assignment}{Cardinal assignment}
\begin{proof}
\emph{Item~1.} Assume the Axiom of Choice and let \(A\) be a set. By
\cref{thm:p01-w02-p01-008-well-ordering}, some \(\prec\) strictly
well-orders \(A\); by \cref{thm:p01-w02-p01-008-representation}, item~1,
\((A,\prec)\cong(\alpha,\in)\) for an ordinal \(\alpha\), so \(A\) is
bijective with \(\alpha\). The set \(\alpha+1\) is an ordinal (ordinal
item~5). By separation let \(K\) be the set of \(\beta\in\alpha+1\) that are
bijective with \(A\). Then \(\alpha\in K\), so \(K\) has a least member
\(\kappa\) (ordinal items~1 and~4). If an ordinal \(\beta\) is bijective
with \(A\) and \(\beta<\kappa\), then \(\beta\in\kappa\subseteq\alpha+1\),
since \(\alpha+1\) is transitive; so \(\beta\in K\) and \(\beta<\kappa\),
against the choice of \(\kappa\). Hence \(\kappa\) is the least ordinal
bijective with \(A\). If \(\gamma<\kappa\) were bijective with \(\kappa\),
composing bijections would make \(\gamma\) bijective with \(A\), which was
just excluded. So \(\kappa\) is a cardinal.

\emph{Item~2.} Let \(\kappa\neq\lambda\) be cardinals. By ordinal item~3,
\(\kappa\in\lambda\) or \(\lambda\in\kappa\); say \(\kappa\in\lambda\). A
bijection between them would make \(\lambda\) bijective with the smaller
ordinal \(\kappa\), against the definition of cardinal. If a set were
bijective with two cardinals, composing bijections would make them
bijective with each other, so they would be equal. With item~1, under the
Axiom of Choice each set is bijective with exactly one cardinal.

\emph{Item~3.} Fix bijections \(p:A\to\kappa\) and \(q:B\to\lambda\). If
\(\kappa\leq\lambda\), then \(\kappa\subseteq\lambda\) (ordinal item~3); with
\(\iota:\kappa\to\lambda\) the inclusion map, \(q^{-1}\circ\iota\circ p\) is an
injection \(A\to B\). Conversely let
\(f:A\to B\) be injective, and suppose \(\kappa\leq\lambda\) fails, so
\(\lambda\in\kappa\) by ordinal item~3. The map
\(g=q\circ f\circ p^{-1}:\kappa\to\lambda\) is injective. Let
\(Y=g[\kappa]\subseteq\lambda\); then \(g\) is a bijection
\(\kappa\to Y\), and \(Y\) is strictly well-ordered by \(\in\) as a subset of
an ordinal. By \cref{thm:p01-w02-p01-008-representation},
\((Y,\in)\cong(\delta,\in)\) for an ordinal \(\delta\) with
\(\delta\leq\lambda\). Since \(\lambda\in\kappa\) and \(\kappa\) is
transitive, \(\delta\in\kappa\) (either \(\delta=\lambda\), or
\(\delta\in\lambda\in\kappa\)). But \(\kappa\) is bijective with \(Y\) and
so with the smaller ordinal \(\delta\), against the definition of
cardinal. Hence an injection \(A\to B\) exists iff \(\kappa\leq\lambda\).

A bijection \(A\to B\) gives the bijection \(q\circ f\circ p^{-1}\) between
\(\kappa\) and \(\lambda\), so \(\kappa=\lambda\) by item~2; conversely if
\(\kappa=\lambda\), \(q^{-1}\circ p\) is a bijection \(A\to B\). Combining
the two equivalences, an injection without a bijection exists iff
\(\kappa\leq\lambda\) and \(\kappa\neq\lambda\), that is, iff
\(\kappa<\lambda\). Under the Axiom of Choice take \(\kappa=\abs A\) and
\(\lambda=\abs B\) from items~1 and~2: the relations \(\abs A\leq\abs B\),
\(\abs A=\abs B\), \(\abs A<\abs B\) of \cref{sec:set-prerequisites-05}
hold iff \(\kappa\leq\lambda\), \(\kappa=\lambda\), \(\kappa<\lambda\)
respectively.

\emph{Item~4.} By ordinal item~3, \(\kappa\in\lambda\), \(\kappa=\lambda\)
or \(\lambda\in\kappa\), so \(\kappa\leq\lambda\) or \(\lambda\leq\kappa\).
Under the Axiom of Choice apply this to \(\abs A\) and \(\abs B\) and use
item~3. For the maximum, recall that for ordinals \(\nu\leq\mu\) iff
\(\nu\subseteq\mu\) (ordinal item~3), so \(\leq\) is transitive on ordinals.
We prove, by induction on \(m\in\N\), that every function \(k\) on \(m+1\)
whose values are ordinals has a greatest value: some \(j\in m+1\) with
\(k(i)\leq k(j)\) for every \(i\in m+1\). For \(m=0\), \(m+1=\{0\}\) and
\(j=0\) serves. Let \(k\) be a function on \((m+1)+1\) with ordinal values.
By the induction hypothesis its restriction to \(m+1\) has a greatest value
\(k(j')\). By ordinal item~3, \(k(j')\leq k(m+1)\) or \(k(m+1)\leq k(j')\);
let \(j\) be \(m+1\) in the first case and \(j'\) in the second. Every
\(i\in(m+1)+1\) lies in \(m+1\) or equals \(m+1\), and in each case
\(k(i)\leq k(j)\) by transitivity of \(\leq\). A nonempty finite list
\(\kappa_1,\ldots,\kappa_n\) of cardinals is a function on a nonzero
natural number, which has the form \(m+1\) by item~4 of
\cref{thm:p01-w02-p01-008-naturals}; its greatest value is
\(\max(\kappa_1,\ldots,\kappa_n)\).

\emph{Item~5.} By \cref{thm:p01-w02-p01-008-pigeonhole}, item~3, \(n\) is a
cardinal; it is bijective with \(n\), so by item~2 it is the only cardinal
bijective with \(n\). By item~4 there, \(\N=\omega\) is a cardinal, and it
is the only cardinal bijective with \(\N\). None of these steps uses the
Axiom of Choice.

\emph{Item~6.} Suppose every set is bijective with some ordinal, and let
\(X\) be a set of nonempty sets. Let \(h:\bigcup X\to\alpha\) be a bijection
onto an ordinal \(\alpha\). For \(B\in X\), \(h[B]\) is a nonempty set of
ordinals, with a least member by ordinal item~4. By separation from
\(X\times\bigcup X\), let \(c\) be the set of pairs \((B,u)\) with \(u\in B\)
and \(h(u)\leq h(v)\) for every \(v\in B\). For each \(B\in X\) such a \(u\)
exists, namely \(h^{-1}\) of the least member of \(h[B]\); and it is unique,
since \(h(u)\leq h(u')\) and \(h(u')\leq h(u)\) give \(h(u)=h(u')\) by
ordinal item~3, hence \(u=u'\). So \(c:X\to\bigcup X\) is a function with
\(c(B)\in B\) for every \(B\in X\), which is the Axiom of Choice for \(X\).
\end{proof}
\end{annexitem}
